% Part 3 -- Bui Duc Phan Hoang: web interface and question answering
% (the live web demo follows slide 16 and belongs to Part 4, Do Hoang Dung)

% --------- Slide 15: Web interface ---------
\begin{frame}{Web Interface: Four Tabs}
    \centering
    \includegraphics[height=5.1cm]{ui_resource.png}\hfill
    \includegraphics[height=5.1cm]{ui_sparql.png}

    \vspace{0.5em}
    \raggedright
    \footnotesize
    \begin{itemize}
        \item \textbf{Resource tree}: 8{,}624 resources under the \term{vio:} class tree
        \item \textbf{Resource} (left): page like \term{dbpedia.org/page/…}; classes marked
              \emph{khai báo} (asserted) or \emph{suy luận} (inferred); graph and relation tree
        \item \textbf{Question answering}: ask in Vietnamese (next slide)
        \item \textbf{SPARQL} (right): editor, 19 prefixes, 9 example queries, table/JSON/CSV
    \end{itemize}
\end{frame}

% --------- Slide 16: Question answering ---------
\begin{frame}{Question Answering: the LLM Writes SPARQL}
    \centering
    \adjustbox{max width=\textwidth}{%
    \begin{tikzpicture}[
        font=\scriptsize, >={Stealth[length=1.8mm]},
        box/.style={draw=black!55, rounded corners=2pt, align=center, text width=2.6cm,
                    minimum height=1.1cm, inner sep=3pt},
        stage/.style={box, draw=hustred, fill=hustlight, line width=0.8pt},
        io/.style={box, fill=black!7},
        flow/.style={->, black!60, thick},
      ]
      \node[io]                            (q)   {Question\\ in Vietnamese};
      \node[stage, right=0.5cm of q]       (pr)  {Prompt: schema\\ summary, rules,\\ 6 examples};
      \node[stage, right=0.5cm of pr]      (llm) {LLM\\ (any OpenAI-\\compatible API)};
      \node[io, right=0.5cm of llm]        (sp)  {SPARQL query};
      \node[stage, below=0.9cm of sp]      (val) {Validate:\\ only \texttt{SELECT}\\ or \texttt{ASK}};
      \node[stage, left=0.5cm of val]      (run) {Run on the\\ rdflib graph};
      \node[stage, left=0.5cm of run]      (ans) {LLM answers\\ from at most\\ 30 rows};
      \node[io, left=0.5cm of ans]         (out) {Answer +\\ reasoning};
      \draw[flow] (q) -- (pr);
      \draw[flow] (pr) -- (llm);
      \draw[flow] (llm) -- (sp);
      \draw[flow] (sp) -- (val);
      \draw[flow] (val) -- (run);
      \draw[flow] (run) -- (ans);
      \draw[flow] (ans) -- (out);
      \draw[->, hustred, dashed, thick] (run.north) -- node[right, font=\tiny, text=hustred, align=left]
            {error or 0 rows:\\ repair once} (llm.south);
    \end{tikzpicture}}

    \vspace{0.8em}
    \raggedright
    \begin{itemize}
        \footnotesize
        \item \textbf{SPARQL mode}: show the query and the raw rows, skip the answer call
        \item \textbf{Entity lookup by name in a subquery}: rdflib applies
              \term{FILTER} after the join, so the Công Phượng career query
              drops from \alert{10.7\,s to 0.27\,s}
        \item Works with OpenAI, Gemini or a local Ollama model (\term{.env});
              tests use a stubbed model, no API key needed
    \end{itemize}
\end{frame}
